Our work draws on several intersecting areas of research.

\paragraph{Abstractive and extractive summarization.}
Text summarization has been a long-standing goal of NLP, with approaches ranging from extractive methods that select key sentences \citep{devlin2019bert} to abstractive methods that generate novel text conditioned on source documents \citep{vaswani2017attention}. BibleGoAI employs a hybrid approach: the quote extraction phase is extractive (selecting verbatim passages), while the BHT generation phase is abstractive (synthesizing a coherent paragraph from selected quotes). This hybrid design allows us to leverage the fluency of LLMs while maintaining explicit grounding in source material.

\paragraph{Retrieval-augmented generation.}
RAG systems \citep{lewis2020rag} address the hallucination problem inherent in LLMs by conditioning generation on retrieved documents. Our system applies RAG specifically to multi-verse biblical commentary where per-verse segmentation is impractical, and organizes the retrieved context into a tiered structure across multiple independent commentaries before generation---a design choice that differentiates it from typical RAG systems that retrieve from a single corpus.

\paragraph{Semantic similarity and sentence embeddings.}
The Universal Sentence Encoder \citep{cer2018use} and Sentence-BERT \citep{reimers2019sentencebert} provide dense vector representations that enable fine-grained semantic comparison. We use these embeddings for three distinct purposes: retrieval indexing (via FAISS \citep{johnson2019faiss}), quality scoring (measuring BHT-to-commentary and BHT-to-verse similarity), and footnote generation (linking BHT sentences to source quotes via cross-encoder scoring).

\paragraph{AI in digital humanities and biblical studies.}
The application of NLP to religious and historical texts is a growing field. Prior work has explored machine translation of biblical texts \citep{nllb2022translate}, named entity recognition in historical documents, and topic modeling across theological corpora. BibleGoAI contributes to this space by focusing specifically on \textit{commentary synthesis}---a task that requires not only understanding the source texts but also respecting the scholarly authority and interpretive nuances of individual commentators.

\paragraph{Existing digital commentary platforms.}
Several widely used platforms make public-domain biblical commentary available online, including BibleHub, BlueLetterBible, StudyLight, and Logos Bible Software (commercial). These systems primarily serve as \emph{browsers} over digitized commentary: a reader selects a verse and sees the corresponding excerpts from each commentator, one commentator at a time. They do not attempt to synthesize across commentators, enforce length budgets, or link generated text back to source quotes. BibleGoAI complements rather than replaces these tools: the source corpus itself is drawn from StudyLight, and BibleGo.org links out to full commentator text for readers who want to dive deeper. The contribution here is the \emph{synthesis layer}---producing a short, source-grounded summary that a reader can consume in seconds rather than a multi-commentator reading session---together with the provenance mechanism that keeps each synthesized claim traceable to the underlying commentary.
